%%%%%%%%%%%%%%%%%%%%%%%%%%%%%%%%%%%%%%%%%%%%%%%%%%%%%%%%%%%%
% 6.5 ENUMERATIVE COMBINATORICS
% The Composition of Advanced Mathematics
%%%%%%%%%%%%%%%%%%%%%%%%%%%%%%%%%%%%%%%%%%%%%%%%%%%%%%%%%%%%

\section{Enumerative Combinatorics}
\label{sec:enumerative-combinatorics}

\prerequisite{
Basic counting, recurrence relations, generating functions, permutations,
partitions, binomial coefficients, and familiarity with combinatorial
classes.
}

\learningobjectives{
By the end of this section, the reader should be able to:
\begin{itemize}
    \item explain the central goals of enumerative combinatorics;
    \item count structured combinatorial classes by size;
    \item distinguish labeled and unlabeled enumeration;
    \item use recurrence relations and generating functions in
          enumeration;
    \item work with Stirling numbers of both kinds;
    \item work with Bell numbers;
    \item count integer partitions;
    \item recognize Catalan families;
    \item count lattice paths;
    \item count permutation statistics such as cycles, inversions, and
          descents;
    \item use Eulerian numbers;
    \item recognize partition identities;
    \item interpret counting formulas bijectively;
    \item use inclusion--exclusion in enumerative problems;
    \item understand the role of asymptotic enumeration;
    \item connect discrete structures through shared counting sequences.
\end{itemize}
}

%%%%%%%%%%%%%%%%%%%%%%%%%%%%%%%%%%%%%%%%%%%%%%%%%%%%%%%%%%%%
\subsection{What Is Enumerative Combinatorics?}
\label{subsec:what-is-enumerative-combinatorics}
%%%%%%%%%%%%%%%%%%%%%%%%%%%%%%%%%%%%%%%%%%%%%%%%%%%%%%%%%%%%

Enumerative combinatorics studies systematic methods for determining
the number of discrete structures satisfying specified conditions.

The basic problem is

\[
\boxed{
c_n
=
|\mathcal{C}_n|,
}
\]

where

\[
\mathcal{C}_n
\]

is the family of structures of size \(n\).

The goal is not merely to count isolated examples.

Instead, one seeks formulas, recurrences, generating functions, and
structural explanations for entire counting sequences.

%%%%%%%%%%%%%%%%%%%%%%%%%%%%%%%%%%%%%%%%%%%%%%%%%%%%%%%%%%%%
\subsection{Counting Sequences}
\label{subsec:enumerative-counting-sequences}
%%%%%%%%%%%%%%%%%%%%%%%%%%%%%%%%%%%%%%%%%%%%%%%%%%%%%%%%%%%%

A combinatorial class

\[
\mathcal{C}
=
\mathcal{C}_0
\cup
\mathcal{C}_1
\cup
\mathcal{C}_2
\cup\cdots
\]

produces a sequence

\[
c_0,c_1,c_2,\ldots,
\]

where

\[
c_n=|\mathcal{C}_n|.
\]

Examples include:

\[
2^n
\]

for subsets of an \(n\)-element set,

\[
n!
\]

for permutations,

and

\[
C_n
=
\frac{1}{n+1}\binom{2n}{n}
\]

for Catalan structures.

%%%%%%%%%%%%%%%%%%%%%%%%%%%%%%%%%%%%%%%%%%%%%%%%%%%%%%%%%%%%
\subsection{Equivalent Combinatorial Classes}
\label{subsec:equivalent-combinatorial-classes}
%%%%%%%%%%%%%%%%%%%%%%%%%%%%%%%%%%%%%%%%%%%%%%%%%%%%%%%%%%%%

Two apparently different combinatorial classes may have the same
counting sequence.

For example, Catalan numbers count:

\begin{itemize}
    \item balanced parenthesis strings;
    \item full binary trees;
    \item triangulations of a convex polygon;
    \item certain lattice paths;
    \item noncrossing partitions;
    \item rooted plane trees.
\end{itemize}

When two families have the same counting sequence, one seeks a
bijection explaining why.

Thus enumerative combinatorics asks not only

\[
\text{``how many?''}
\]

but also

\[
\boxed{
\text{``why do these different objects have the same count?''}
}
\]

%%%%%%%%%%%%%%%%%%%%%%%%%%%%%%%%%%%%%%%%%%%%%%%%%%%%%%%%%%%%
\subsection{Labeled and Unlabeled Enumeration}
\label{subsec:labeled-unlabeled-enumeration}
%%%%%%%%%%%%%%%%%%%%%%%%%%%%%%%%%%%%%%%%%%%%%%%%%%%%%%%%%%%%

A labeled object has distinguishable components.

An unlabeled object is considered only up to structural equivalence.

For example, on a labeled vertex set

\[
\{1,2,\ldots,n\},
\]

the number of simple graphs is

\[
\boxed{
2^{\binom{n}{2}}.
}
\]

Counting unlabeled graphs is substantially harder because graphs that
differ only by relabeling are identified.

This distinction is fundamental throughout enumeration.

%%%%%%%%%%%%%%%%%%%%%%%%%%%%%%%%%%%%%%%%%%%%%%%%%%%%%%%%%%%%
\subsection{Permutations as Enumerative Objects}
\label{subsec:enumerating-permutations}
%%%%%%%%%%%%%%%%%%%%%%%%%%%%%%%%%%%%%%%%%%%%%%%%%%%%%%%%%%%%

The symmetric group

\[
S_n
\]

contains all permutations of \(n\) elements.

Ignoring its algebraic group structure for the moment,

\[
\boxed{
|S_n|=n!.
}
\]

Enumerative combinatorics refines this total count by classifying
permutations according to statistics such as:

\begin{itemize}
    \item number of cycles;
    \item number of fixed points;
    \item number of inversions;
    \item number of descents;
    \item cycle type.
\end{itemize}

%%%%%%%%%%%%%%%%%%%%%%%%%%%%%%%%%%%%%%%%%%%%%%%%%%%%%%%%%%%%
\subsection{Cycle Decomposition}
\label{subsec:permutation-cycle-decomposition}
%%%%%%%%%%%%%%%%%%%%%%%%%%%%%%%%%%%%%%%%%%%%%%%%%%%%%%%%%%%%

Every permutation decomposes uniquely into disjoint cycles, up to the
order in which the cycles are written and cyclic rotation within each
cycle.

For example,

\[
\pi
=
(1\,4\,3)(2\,5)
\]

has two cycles.

Cycle decompositions lead naturally to Stirling numbers of the first
kind.

%%%%%%%%%%%%%%%%%%%%%%%%%%%%%%%%%%%%%%%%%%%%%%%%%%%%%%%%%%%%
\subsection{Stirling Numbers of the First Kind}
\label{subsec:stirling-first-kind}
%%%%%%%%%%%%%%%%%%%%%%%%%%%%%%%%%%%%%%%%%%%%%%%%%%%%%%%%%%%%

\begin{definitionbox}{Unsigned Stirling Number of the First Kind}
The unsigned Stirling number of the first kind

\[
\boxed{
\left[
\begin{matrix}
n\\
k
\end{matrix}
\right]
}
\]

counts permutations of \(n\) elements having exactly \(k\) disjoint
cycles.
\end{definitionbox}

The notation is read ``Stirling \(n\) over \(k\), first kind.''

For example,

\[
\left[
\begin{matrix}
3\\
1
\end{matrix}
\right]
=2,
\]

because the permutations consisting of one \(3\)-cycle are

\[
(1\,2\,3)
\]

and

\[
(1\,3\,2).
\]

%%%%%%%%%%%%%%%%%%%%%%%%%%%%%%%%%%%%%%%%%%%%%%%%%%%%%%%%%%%%
\subsection{Recurrence for First-Kind Stirling Numbers}
\label{subsec:stirling-first-recurrence}
%%%%%%%%%%%%%%%%%%%%%%%%%%%%%%%%%%%%%%%%%%%%%%%%%%%%%%%%%%%%

The recurrence is

\[
\boxed{
\left[
\begin{matrix}
n\\
k
\end{matrix}
\right]
=
\left[
\begin{matrix}
n-1\\
k-1
\end{matrix}
\right]
+
(n-1)
\left[
\begin{matrix}
n-1\\
k
\end{matrix}
\right].
}
\]

To understand this, insert element \(n\).

Either:

\begin{itemize}
    \item \(n\) forms a new cycle by itself, or
    \item \(n\) is inserted into one of the existing cycles.
\end{itemize}

If there are \(n-1\) elements already arranged in cycles, there are
\(n-1\) possible insertion positions across all cycles.

%%%%%%%%%%%%%%%%%%%%%%%%%%%%%%%%%%%%%%%%%%%%%%%%%%%%%%%%%%%%
\subsection{Worked Example: First-Kind Stirling Numbers}
\label{subsec:worked-stirling-first}
%%%%%%%%%%%%%%%%%%%%%%%%%%%%%%%%%%%%%%%%%%%%%%%%%%%%%%%%%%%%

\begin{workedexamplebox}{Permutations of Four Elements with Two Cycles}

Use the recurrence

\[
\left[
\begin{matrix}
4\\
2
\end{matrix}
\right]
=
\left[
\begin{matrix}
3\\
1
\end{matrix}
\right]
+
3
\left[
\begin{matrix}
3\\
2
\end{matrix}
\right].
\]

We know

\[
\left[
\begin{matrix}
3\\
1
\end{matrix}
\right]
=2
\]

and

\[
\left[
\begin{matrix}
3\\
2
\end{matrix}
\right]
=3.
\]

Therefore,

\[
2+3(3)=11.
\]

\begin{answerbox}
\[
\boxed{
\left[
\begin{matrix}
4\\
2
\end{matrix}
\right]
=
11.
}
\]
\end{answerbox}

\end{workedexamplebox}

%%%%%%%%%%%%%%%%%%%%%%%%%%%%%%%%%%%%%%%%%%%%%%%%%%%%%%%%%%%%
\subsection{Signed Stirling Numbers of the First Kind}
\label{subsec:signed-stirling-first}
%%%%%%%%%%%%%%%%%%%%%%%%%%%%%%%%%%%%%%%%%%%%%%%%%%%%%%%%%%%%

A signed version is often defined by

\[
\boxed{
s(n,k)
=
(-1)^{n-k}
\left[
\begin{matrix}
n\\
k
\end{matrix}
\right].
}
\]

These numbers appear in polynomial expansions involving falling
factorials.

For example,

\[
(x)_n
=
x(x-1)\cdots(x-n+1)
\]

can be expanded as

\[
\boxed{
(x)_n
=
\sum_{k=0}^{n}
s(n,k)x^k.
}
\]

%%%%%%%%%%%%%%%%%%%%%%%%%%%%%%%%%%%%%%%%%%%%%%%%%%%%%%%%%%%%
\subsection{Stirling Numbers of the Second Kind}
\label{subsec:stirling-second-kind}
%%%%%%%%%%%%%%%%%%%%%%%%%%%%%%%%%%%%%%%%%%%%%%%%%%%%%%%%%%%%

\begin{definitionbox}{Stirling Number of the Second Kind}
The Stirling number of the second kind

\[
\boxed{
\left\{
\begin{matrix}
n\\
k
\end{matrix}
\right\}
}
\]

counts partitions of an \(n\)-element set into exactly \(k\) nonempty
unlabeled blocks.
\end{definitionbox}

For example,

\[
\left\{
\begin{matrix}
3\\
2
\end{matrix}
\right\}
=3.
\]

The partitions are

\[
\{\{1\},\{2,3\}\},
\]

\[
\{\{2\},\{1,3\}\},
\]

and

\[
\{\{3\},\{1,2\}\}.
\]

%%%%%%%%%%%%%%%%%%%%%%%%%%%%%%%%%%%%%%%%%%%%%%%%%%%%%%%%%%%%
\subsection{Recurrence for Second-Kind Stirling Numbers}
\label{subsec:stirling-second-recurrence}
%%%%%%%%%%%%%%%%%%%%%%%%%%%%%%%%%%%%%%%%%%%%%%%%%%%%%%%%%%%%

The recurrence is

\[
\boxed{
\left\{
\begin{matrix}
n\\
k
\end{matrix}
\right\}
=
\left\{
\begin{matrix}
n-1\\
k-1
\end{matrix}
\right\}
+
k
\left\{
\begin{matrix}
n-1\\
k
\end{matrix}
\right\}.
}
\]

Consider the element \(n\).

Either:

\begin{itemize}
    \item \(n\) forms a new singleton block, or
    \item \(n\) joins one of the \(k\) existing blocks.
\end{itemize}

These cases are disjoint and exhaustive.

%%%%%%%%%%%%%%%%%%%%%%%%%%%%%%%%%%%%%%%%%%%%%%%%%%%%%%%%%%%%
\subsection{Worked Example: Second-Kind Stirling Numbers}
\label{subsec:worked-stirling-second}
%%%%%%%%%%%%%%%%%%%%%%%%%%%%%%%%%%%%%%%%%%%%%%%%%%%%%%%%%%%%

\begin{workedexamplebox}{Partitioning Five Elements into Two Blocks}

Use

\[
\left\{
\begin{matrix}
5\\
2
\end{matrix}
\right\}
=
\left\{
\begin{matrix}
4\\
1
\end{matrix}
\right\}
+
2
\left\{
\begin{matrix}
4\\
2
\end{matrix}
\right\}.
\]

Since

\[
\left\{
\begin{matrix}
4\\
1
\end{matrix}
\right\}
=1
\]

and

\[
\left\{
\begin{matrix}
4\\
2
\end{matrix}
\right\}
=7,
\]

we obtain

\[
1+2(7)=15.
\]

\begin{answerbox}
\[
\boxed{
\left\{
\begin{matrix}
5\\
2
\end{matrix}
\right\}
=
15.
}
\]
\end{answerbox}

\end{workedexamplebox}

%%%%%%%%%%%%%%%%%%%%%%%%%%%%%%%%%%%%%%%%%%%%%%%%%%%%%%%%%%%%
\subsection{Explicit Formula for Second-Kind Stirling Numbers}
\label{subsec:stirling-second-explicit}
%%%%%%%%%%%%%%%%%%%%%%%%%%%%%%%%%%%%%%%%%%%%%%%%%%%%%%%%%%%%

Using inclusion--exclusion,

\[
\boxed{
\left\{
\begin{matrix}
n\\
k
\end{matrix}
\right\}
=
\frac{1}{k!}
\sum_{j=0}^{k}
(-1)^{k-j}
\binom{k}{j}
j^n.
}
\]

The factor \(k!\) appears because a set partition has unlabeled blocks,
while a surjective function onto \(k\) labeled targets labels the
blocks.

%%%%%%%%%%%%%%%%%%%%%%%%%%%%%%%%%%%%%%%%%%%%%%%%%%%%%%%%%%%%
\subsection{Bell Numbers}
\label{subsec:bell-numbers}
%%%%%%%%%%%%%%%%%%%%%%%%%%%%%%%%%%%%%%%%%%%%%%%%%%%%%%%%%%%%

\begin{definitionbox}{Bell Number}
The Bell number \(B_n\) counts all partitions of an \(n\)-element set.
\end{definitionbox}

Since a partition may contain any number of blocks,

\[
\boxed{
B_n
=
\sum_{k=0}^{n}
\left\{
\begin{matrix}
n\\
k
\end{matrix}
\right\}.
}
\]

The first Bell numbers are

\[
1,1,2,5,15,52,203,\ldots.
\]

%%%%%%%%%%%%%%%%%%%%%%%%%%%%%%%%%%%%%%%%%%%%%%%%%%%%%%%%%%%%
\subsection{Bell Number Recurrence}
\label{subsec:bell-number-recurrence}
%%%%%%%%%%%%%%%%%%%%%%%%%%%%%%%%%%%%%%%%%%%%%%%%%%%%%%%%%%%%

A useful recurrence is

\[
\boxed{
B_{n+1}
=
\sum_{k=0}^{n}
\binom{n}{k}B_k.
}
\]

To see this, focus on the block containing a distinguished element.

Choose \(k\) of the remaining \(n\) elements to lie outside that block,
then partition those \(k\) elements arbitrarily.

%%%%%%%%%%%%%%%%%%%%%%%%%%%%%%%%%%%%%%%%%%%%%%%%%%%%%%%%%%%%
\subsection{Bell Number Exponential Generating Function}
\label{subsec:bell-number-egf}
%%%%%%%%%%%%%%%%%%%%%%%%%%%%%%%%%%%%%%%%%%%%%%%%%%%%%%%%%%%%

The exponential generating function is

\[
\boxed{
\sum_{n=0}^{\infty}
B_n\frac{x^n}{n!}
=
e^{e^x-1}.
}
\]

This compact formula reflects the fact that a set partition is a set of
nonempty labeled blocks.

%%%%%%%%%%%%%%%%%%%%%%%%%%%%%%%%%%%%%%%%%%%%%%%%%%%%%%%%%%%%
\subsection{Integer Partitions}
\label{subsec:enumerative-integer-partitions}
%%%%%%%%%%%%%%%%%%%%%%%%%%%%%%%%%%%%%%%%%%%%%%%%%%%%%%%%%%%%

An integer partition of \(n\) is a representation

\[
n
=
\lambda_1+\lambda_2+\cdots+\lambda_k
\]

with

\[
\lambda_1\geq\lambda_2\geq\cdots\geq\lambda_k\geq1.
\]

The Greek letter

\[
\lambda
\]

is called ``lambda.''

The number of integer partitions of \(n\) is denoted

\[
\boxed{
p(n).
}
\]

%%%%%%%%%%%%%%%%%%%%%%%%%%%%%%%%%%%%%%%%%%%%%%%%%%%%%%%%%%%%
\subsection{Worked Example: Partitions of Five}
\label{subsec:worked-partitions-five}
%%%%%%%%%%%%%%%%%%%%%%%%%%%%%%%%%%%%%%%%%%%%%%%%%%%%%%%%%%%%

\begin{workedexamplebox}{Integer Partitions of \(5\)}

The partitions of \(5\) are

\[
5,
\]

\[
4+1,
\]

\[
3+2,
\]

\[
3+1+1,
\]

\[
2+2+1,
\]

\[
2+1+1+1,
\]

\[
1+1+1+1+1.
\]

\begin{answerbox}
\[
\boxed{
p(5)=7.
}
\]
\end{answerbox}

\end{workedexamplebox}

%%%%%%%%%%%%%%%%%%%%%%%%%%%%%%%%%%%%%%%%%%%%%%%%%%%%%%%%%%%%
\subsection{Partition Generating Function}
\label{subsec:partition-generating-function-enumerative}
%%%%%%%%%%%%%%%%%%%%%%%%%%%%%%%%%%%%%%%%%%%%%%%%%%%%%%%%%%%%

Euler's generating function is

\[
\boxed{
\sum_{n=0}^{\infty}p(n)x^n
=
\prod_{m=1}^{\infty}\frac{1}{1-x^m}.
}
\]

The factor

\[
\frac{1}{1-x^m}
\]

records the possible multiplicities of a part of size \(m\).

%%%%%%%%%%%%%%%%%%%%%%%%%%%%%%%%%%%%%%%%%%%%%%%%%%%%%%%%%%%%
\subsection{Ferrers Diagrams}
\label{subsec:ferrers-diagrams}
%%%%%%%%%%%%%%%%%%%%%%%%%%%%%%%%%%%%%%%%%%%%%%%%%%%%%%%%%%%%

An integer partition may be represented using rows of dots.

For example,

\[
7=4+2+1
\]

has Ferrers diagram

\[
\begin{array}{cccc}
\bullet & \bullet & \bullet & \bullet\\
\bullet & \bullet\\
\bullet
\end{array}
\]

Ferrers diagrams convert partition identities into geometric
transformations.

%%%%%%%%%%%%%%%%%%%%%%%%%%%%%%%%%%%%%%%%%%%%%%%%%%%%%%%%%%%%
\subsection{Conjugate Partitions}
\label{subsec:conjugate-partitions}
%%%%%%%%%%%%%%%%%%%%%%%%%%%%%%%%%%%%%%%%%%%%%%%%%%%%%%%%%%%%

Reflecting a Ferrers diagram across its main diagonal produces the
conjugate partition.

This establishes a bijection between:

\[
\boxed{
\text{partitions of }n\text{ with at most }k\text{ parts}
}
\]

and

\[
\boxed{
\text{partitions of }n\text{ with largest part at most }k.
}
\]

This is a classic example of a bijective proof.

%%%%%%%%%%%%%%%%%%%%%%%%%%%%%%%%%%%%%%%%%%%%%%%%%%%%%%%%%%%%
\subsection{Partitions into Distinct Parts}
\label{subsec:enumerative-distinct-partitions}
%%%%%%%%%%%%%%%%%%%%%%%%%%%%%%%%%%%%%%%%%%%%%%%%%%%%%%%%%%%%

A partition into distinct parts uses each part size at most once.

Its generating function is

\[
\boxed{
\prod_{m=1}^{\infty}(1+x^m).
}
\]

The term \(1\) represents not using part \(m\).

The term \(x^m\) represents using it once.

%%%%%%%%%%%%%%%%%%%%%%%%%%%%%%%%%%%%%%%%%%%%%%%%%%%%%%%%%%%%
\subsection{Partitions into Odd Parts}
\label{subsec:enumerative-odd-partitions}
%%%%%%%%%%%%%%%%%%%%%%%%%%%%%%%%%%%%%%%%%%%%%%%%%%%%%%%%%%%%

A partition into odd parts uses only

\[
1,3,5,\ldots.
\]

Its generating function is

\[
\boxed{
\prod_{m=1}^{\infty}
\frac{1}{1-x^{2m-1}}.
}
\]

Euler proved the identity

\[
\boxed{
\prod_{m=1}^{\infty}(1+x^m)
=
\prod_{m=1}^{\infty}
\frac{1}{1-x^{2m-1}}.
}
\]

Therefore,

\[
\boxed{
\text{partitions into distinct parts}
=
\text{partitions into odd parts}.
}
\]

%%%%%%%%%%%%%%%%%%%%%%%%%%%%%%%%%%%%%%%%%%%%%%%%%%%%%%%%%%%%
\subsection{Catalan Numbers}
\label{subsec:catalan-numbers-enumerative}
%%%%%%%%%%%%%%%%%%%%%%%%%%%%%%%%%%%%%%%%%%%%%%%%%%%%%%%%%%%%

The Catalan numbers are

\[
\boxed{
C_n
=
\frac{1}{n+1}
\binom{2n}{n}.
}
\]

The first values are

\[
1,1,2,5,14,42,132,\ldots.
\]

They satisfy the recurrence

\[
\boxed{
C_{n+1}
=
\sum_{k=0}^{n}C_kC_{n-k}.
}
\]

Their generating function is

\[
\boxed{
C(x)
=
\frac{1-\sqrt{1-4x}}{2x}.
}
\]

%%%%%%%%%%%%%%%%%%%%%%%%%%%%%%%%%%%%%%%%%%%%%%%%%%%%%%%%%%%%
\subsection{Dyck Paths}
\label{subsec:dyck-paths}
%%%%%%%%%%%%%%%%%%%%%%%%%%%%%%%%%%%%%%%%%%%%%%%%%%%%%%%%%%%%

A Dyck path of semilength \(n\) is a lattice path from

\[
(0,0)
\]

to

\[
(2n,0)
\]

using steps

\[
U=(1,1)
\]

and

\[
D=(1,-1)
\]

that never goes below the horizontal axis.

The number of Dyck paths of semilength \(n\) is

\[
\boxed{
C_n.
}
\]

%%%%%%%%%%%%%%%%%%%%%%%%%%%%%%%%%%%%%%%%%%%%%%%%%%%%%%%%%%%%
\subsection{Balanced Parentheses}
\label{subsec:balanced-parentheses}
%%%%%%%%%%%%%%%%%%%%%%%%%%%%%%%%%%%%%%%%%%%%%%%%%%%%%%%%%%%%

A balanced parenthesis string with \(n\) pairs of parentheses is another
Catalan object.

For \(n=3\), the five strings are

\[
((())),
\]

\[
(()()),
\]

\[
(())(),
\]

\[
()(()),
\]

\[
()()().
\]

Therefore,

\[
\boxed{
C_3=5.
}
\]

Mapping

\[
(
\longleftrightarrow U
\]

and

\[
)
\longleftrightarrow D
\]

gives a bijection with Dyck paths.

%%%%%%%%%%%%%%%%%%%%%%%%%%%%%%%%%%%%%%%%%%%%%%%%%%%%%%%%%%%%
\subsection{Binary Trees and Catalan Numbers}
\label{subsec:binary-trees-catalan}
%%%%%%%%%%%%%%%%%%%%%%%%%%%%%%%%%%%%%%%%%%%%%%%%%%%%%%%%%%%%

The number of full rooted binary trees with \(n\) internal nodes is

\[
\boxed{
C_n.
}
\]

Decomposing such a tree at its root gives a left subtree and a right
subtree.

If the left subtree has \(k\) internal nodes, the right subtree has

\[
n-1-k.
\]

Therefore,

\[
C_n
=
\sum_{k=0}^{n-1}
C_kC_{n-1-k}.
\]

This is the Catalan recurrence.

%%%%%%%%%%%%%%%%%%%%%%%%%%%%%%%%%%%%%%%%%%%%%%%%%%%%%%%%%%%%
\subsection{Polygon Triangulations}
\label{subsec:polygon-triangulations}
%%%%%%%%%%%%%%%%%%%%%%%%%%%%%%%%%%%%%%%%%%%%%%%%%%%%%%%%%%%%

The number of triangulations of a convex polygon with \(n+2\) vertices
is

\[
\boxed{
C_n.
}
\]

For example, a pentagon has

\[
C_3=5
\]

triangulations.

Thus polygon geometry produces the same sequence as Dyck paths,
parentheses, and binary trees.

%%%%%%%%%%%%%%%%%%%%%%%%%%%%%%%%%%%%%%%%%%%%%%%%%%%%%%%%%%%%
\subsection{Lattice Path Enumeration}
\label{subsec:lattice-path-enumeration}
%%%%%%%%%%%%%%%%%%%%%%%%%%%%%%%%%%%%%%%%%%%%%%%%%%%%%%%%%%%%

A monotone lattice path from

\[
(0,0)
\]

to

\[
(m,n)
\]

using right and upward steps consists of

\[
m+n
\]

total steps.

Choose which \(m\) positions are right steps.

Therefore,

\[
\boxed{
\binom{m+n}{m}
}
\]

such paths exist.

%%%%%%%%%%%%%%%%%%%%%%%%%%%%%%%%%%%%%%%%%%%%%%%%%%%%%%%%%%%%
\subsection{Paths Avoiding a Boundary}
\label{subsec:lattice-path-boundaries}
%%%%%%%%%%%%%%%%%%%%%%%%%%%%%%%%%%%%%%%%%%%%%%%%%%%%%%%%%%%%

More interesting enumeration occurs when a path must avoid a forbidden
region.

For example, paths from

\[
(0,0)
\]

to

\[
(n,n)
\]

that never rise above the diagonal are counted by

\[
\boxed{
C_n
=
\frac{1}{n+1}\binom{2n}{n}.
}
\]

This may be proved using the reflection principle.

%%%%%%%%%%%%%%%%%%%%%%%%%%%%%%%%%%%%%%%%%%%%%%%%%%%%%%%%%%%%
\subsection{Reflection Principle}
\label{subsec:reflection-principle}
%%%%%%%%%%%%%%%%%%%%%%%%%%%%%%%%%%%%%%%%%%%%%%%%%%%%%%%%%%%%

The reflection principle transforms a path crossing a boundary into a
different unrestricted path.

In Catalan counting, start with all

\[
\binom{2n}{n}
\]

monotone paths.

Subtract those that cross the forbidden boundary.

A reflection bijection converts the forbidden paths into another class
that can be counted directly.

The result is

\[
\boxed{
\binom{2n}{n}
-
\binom{2n}{n+1}
=
\frac{1}{n+1}\binom{2n}{n}.
}
\]

%%%%%%%%%%%%%%%%%%%%%%%%%%%%%%%%%%%%%%%%%%%%%%%%%%%%%%%%%%%%
\subsection{Central Binomial Coefficients}
\label{subsec:central-binomial-coefficients}
%%%%%%%%%%%%%%%%%%%%%%%%%%%%%%%%%%%%%%%%%%%%%%%%%%%%%%%%%%%%

The numbers

\[
\boxed{
\binom{2n}{n}
}
\]

are called central binomial coefficients.

They count, among many things:

\begin{itemize}
    \item binary strings with \(n\) zeros and \(n\) ones;
    \item lattice paths from \((0,0)\) to \((n,n)\);
    \item \(n\)-element subsets of a \(2n\)-element set.
\end{itemize}

Their asymptotic behavior is

\[
\boxed{
\binom{2n}{n}
\sim
\frac{4^n}{\sqrt{\pi n}}.
}
\]

This follows from Stirling's approximation.

%%%%%%%%%%%%%%%%%%%%%%%%%%%%%%%%%%%%%%%%%%%%%%%%%%%%%%%%%%%%
\subsection{Stirling's Approximation}
\label{subsec:stirlings-approximation}
%%%%%%%%%%%%%%%%%%%%%%%%%%%%%%%%%%%%%%%%%%%%%%%%%%%%%%%%%%%%

For large \(n\),

\[
\boxed{
n!
\sim
\sqrt{2\pi n}
\left(\frac{n}{e}\right)^n.
}
\]

This is called Stirling's approximation.

It allows exact combinatorial formulas involving factorials to be
converted into asymptotic estimates.

For example,

\[
\binom{2n}{n}
=
\frac{(2n)!}{(n!)^2}
\]

can be simplified asymptotically using this formula.

%%%%%%%%%%%%%%%%%%%%%%%%%%%%%%%%%%%%%%%%%%%%%%%%%%%%%%%%%%%%
\subsection{Asymptotics of Catalan Numbers}
\label{subsec:catalan-asymptotics}
%%%%%%%%%%%%%%%%%%%%%%%%%%%%%%%%%%%%%%%%%%%%%%%%%%%%%%%%%%%%

Since

\[
C_n
=
\frac{1}{n+1}\binom{2n}{n},
\]

and

\[
\binom{2n}{n}
\sim
\frac{4^n}{\sqrt{\pi n}},
\]

we obtain

\[
\boxed{
C_n
\sim
\frac{4^n}{\sqrt{\pi}n^{3/2}}.
}
\]

Thus Catalan structures grow exponentially, with an additional
polynomial correction.

%%%%%%%%%%%%%%%%%%%%%%%%%%%%%%%%%%%%%%%%%%%%%%%%%%%%%%%%%%%%
\subsection{Permutation Inversions}
\label{subsec:permutation-inversions}
%%%%%%%%%%%%%%%%%%%%%%%%%%%%%%%%%%%%%%%%%%%%%%%%%%%%%%%%%%%%

For a permutation

\[
\pi\in S_n,
\]

an inversion is a pair

\[
(i,j)
\]

with

\[
i<j
\]

but

\[
\pi(i)>\pi(j).
\]

The inversion number is denoted

\[
\operatorname{inv}(\pi).
\]

For example,

\[
\pi=312
\]

has inversions

\[
(1,2)
\]

and

\[
(1,3),
\]

so

\[
\operatorname{inv}(\pi)=2.
\]

%%%%%%%%%%%%%%%%%%%%%%%%%%%%%%%%%%%%%%%%%%%%%%%%%%%%%%%%%%%%
\subsection{Mahonian Numbers}
\label{subsec:mahonian-numbers}
%%%%%%%%%%%%%%%%%%%%%%%%%%%%%%%%%%%%%%%%%%%%%%%%%%%%%%%%%%%%

Let

\[
M(n,k)
\]

denote the number of permutations of \(n\) elements having exactly
\(k\) inversions.

These are called Mahonian numbers.

Their generating polynomial is

\[
\boxed{
\sum_{\pi\in S_n}
q^{\operatorname{inv}(\pi)}
=
\prod_{j=1}^{n}
(1+q+\cdots+q^{j-1}).
}
\]

The variable \(q\) tracks inversion number.

%%%%%%%%%%%%%%%%%%%%%%%%%%%%%%%%%%%%%%%%%%%%%%%%%%%%%%%%%%%%
\subsection{Permutation Descents}
\label{subsec:permutation-descents}
%%%%%%%%%%%%%%%%%%%%%%%%%%%%%%%%%%%%%%%%%%%%%%%%%%%%%%%%%%%%

A descent of a permutation \(\pi\in S_n\) is an index \(i\) such that

\[
\pi(i)>\pi(i+1).
\]

The number of descents is written

\[
\operatorname{des}(\pi).
\]

For example,

\[
\pi=3142
\]

has descents at positions \(1\) and \(3\).

Therefore,

\[
\operatorname{des}(\pi)=2.
\]

%%%%%%%%%%%%%%%%%%%%%%%%%%%%%%%%%%%%%%%%%%%%%%%%%%%%%%%%%%%%
\subsection{Eulerian Numbers}
\label{subsec:eulerian-numbers}
%%%%%%%%%%%%%%%%%%%%%%%%%%%%%%%%%%%%%%%%%%%%%%%%%%%%%%%%%%%%

\begin{definitionbox}{Eulerian Number}
The Eulerian number

\[
\boxed{
\left\langle
\begin{matrix}
n\\
k
\end{matrix}
\right\rangle
}
\]

counts permutations of \(n\) elements having exactly \(k\) descents.
\end{definitionbox}

The notation is often read ``Eulerian \(n\) over \(k\).''

These numbers satisfy the recurrence

\[
\boxed{
\left\langle
\begin{matrix}
n\\
k
\end{matrix}
\right\rangle
=
(k+1)
\left\langle
\begin{matrix}
n-1\\
k
\end{matrix}
\right\rangle
+
(n-k)
\left\langle
\begin{matrix}
n-1\\
k-1
\end{matrix}
\right\rangle.
}
\]

%%%%%%%%%%%%%%%%%%%%%%%%%%%%%%%%%%%%%%%%%%%%%%%%%%%%%%%%%%%%
\subsection{Eulerian Polynomials}
\label{subsec:eulerian-polynomials}
%%%%%%%%%%%%%%%%%%%%%%%%%%%%%%%%%%%%%%%%%%%%%%%%%%%%%%%%%%%%

The Eulerian polynomial is

\[
\boxed{
A_n(t)
=
\sum_{\pi\in S_n}
t^{\operatorname{des}(\pi)}.
}
\]

Equivalently,

\[
A_n(t)
=
\sum_{k=0}^{n-1}
\left\langle
\begin{matrix}
n\\
k
\end{matrix}
\right\rangle
t^k.
\]

This is an example of a generating polynomial for a statistic rather
than for size.

%%%%%%%%%%%%%%%%%%%%%%%%%%%%%%%%%%%%%%%%%%%%%%%%%%%%%%%%%%%%
\subsection{Fixed Points of Permutations}
\label{subsec:fixed-points-permutations}
%%%%%%%%%%%%%%%%%%%%%%%%%%%%%%%%%%%%%%%%%%%%%%%%%%%%%%%%%%%%

A fixed point of a permutation is an index \(i\) satisfying

\[
\pi(i)=i.
\]

Derangements have no fixed points.

More generally, the number of permutations of \(n\) elements with
exactly \(k\) fixed points is

\[
\boxed{
\binom{n}{k}D_{n-k}.
}
\]

Choose the \(k\) fixed elements, then derange the remaining \(n-k\)
elements.

%%%%%%%%%%%%%%%%%%%%%%%%%%%%%%%%%%%%%%%%%%%%%%%%%%%%%%%%%%%%
\subsection{Rencontres Numbers}
\label{subsec:rencontres-numbers}
%%%%%%%%%%%%%%%%%%%%%%%%%%%%%%%%%%%%%%%%%%%%%%%%%%%%%%%%%%%%

The numbers

\[
\boxed{
\binom{n}{k}D_{n-k}
}
\]

are often called rencontres numbers.

They refine the total permutation count according to number of fixed
points.

Summing over all possible \(k\),

\[
\boxed{
n!
=
\sum_{k=0}^{n}
\binom{n}{k}D_{n-k}.
}
\]

This partitions permutations by their number of fixed points.

%%%%%%%%%%%%%%%%%%%%%%%%%%%%%%%%%%%%%%%%%%%%%%%%%%%%%%%%%%%%
\subsection{Compositions of Integers}
\label{subsec:integer-compositions}
%%%%%%%%%%%%%%%%%%%%%%%%%%%%%%%%%%%%%%%%%%%%%%%%%%%%%%%%%%%%

A composition of \(n\) is an ordered sequence of positive integers
summing to \(n\).

For example, the compositions of \(4\) include

\[
4,
\]

\[
1+3,
\]

\[
3+1,
\]

\[
2+2,
\]

and others.

Unlike integer partitions, order matters.

%%%%%%%%%%%%%%%%%%%%%%%%%%%%%%%%%%%%%%%%%%%%%%%%%%%%%%%%%%%%
\subsection{Counting Compositions}
\label{subsec:counting-compositions}
%%%%%%%%%%%%%%%%%%%%%%%%%%%%%%%%%%%%%%%%%%%%%%%%%%%%%%%%%%%%

Write

\[
n
=
1+1+\cdots+1
\]

using \(n\) ones.

There are

\[
n-1
\]

gaps between adjacent ones.

In each gap, either insert a separator or do not.

Therefore,

\[
\boxed{
2^{n-1}
}
\]

compositions of \(n\) exist for \(n\geq1\).

%%%%%%%%%%%%%%%%%%%%%%%%%%%%%%%%%%%%%%%%%%%%%%%%%%%%%%%%%%%%
\subsection{Compositions with a Fixed Number of Parts}
\label{subsec:compositions-fixed-parts}
%%%%%%%%%%%%%%%%%%%%%%%%%%%%%%%%%%%%%%%%%%%%%%%%%%%%%%%%%%%%

A composition of \(n\) into exactly \(k\) positive parts corresponds to
a positive integer solution of

\[
x_1+\cdots+x_k=n.
\]

By stars and bars,

\[
\boxed{
\binom{n-1}{k-1}.
}
\]

Summing over \(k\),

\[
\sum_{k=1}^{n}
\binom{n-1}{k-1}
=
2^{n-1}.
\]

%%%%%%%%%%%%%%%%%%%%%%%%%%%%%%%%%%%%%%%%%%%%%%%%%%%%%%%%%%%%
\subsection{Set Partitions Versus Integer Partitions}
\label{subsec:set-versus-integer-partitions}
%%%%%%%%%%%%%%%%%%%%%%%%%%%%%%%%%%%%%%%%%%%%%%%%%%%%%%%%%%%%

These two concepts are different.

A set partition divides a set into blocks.

For example,

\[
\{\{1,4\},\{2\},\{3,5\}\}.
\]

An integer partition expresses a number as an unordered sum.

For example,

\[
5=2+2+1.
\]

Set partitions are counted by Bell and Stirling numbers.

Integer partitions are counted by \(p(n)\).

%%%%%%%%%%%%%%%%%%%%%%%%%%%%%%%%%%%%%%%%%%%%%%%%%%%%%%%%%%%%
\subsection{Compositions Versus Partitions}
\label{subsec:compositions-versus-partitions}
%%%%%%%%%%%%%%%%%%%%%%%%%%%%%%%%%%%%%%%%%%%%%%%%%%%%%%%%%%%%

For compositions, order matters.

Thus,

\[
3+1
\]

and

\[
1+3
\]

are different.

For integer partitions, they are the same.

Therefore,

\[
\boxed{
\text{composition}
=
\text{ordered sum},
}
\]

while

\[
\boxed{
\text{partition}
=
\text{unordered sum}.
}
\]

%%%%%%%%%%%%%%%%%%%%%%%%%%%%%%%%%%%%%%%%%%%%%%%%%%%%%%%%%%%%
\subsection{Binomial Convolution}
\label{subsec:binomial-convolution}
%%%%%%%%%%%%%%%%%%%%%%%%%%%%%%%%%%%%%%%%%%%%%%%%%%%%%%%%%%%%

Many enumerative identities involve sums such as

\[
\sum_{k}
\binom{n}{k}a_kb_{n-k}.
\]

This is a labeled convolution and corresponds naturally to products of
exponential generating functions.

Such expressions arise whenever a labeled set of size \(n\) is divided
between two structures.

%%%%%%%%%%%%%%%%%%%%%%%%%%%%%%%%%%%%%%%%%%%%%%%%%%%%%%%%%%%%
\subsection{Vandermonde's Identity Revisited}
\label{subsec:vandermonde-enumerative}
%%%%%%%%%%%%%%%%%%%%%%%%%%%%%%%%%%%%%%%%%%%%%%%%%%%%%%%%%%%%

Recall

\[
\boxed{
\sum_{k}
\binom{m}{k}
\binom{n}{r-k}
=
\binom{m+n}{r}.
}
\]

This identity can be viewed as an enumerative decomposition.

Select \(r\) objects from two disjoint populations of sizes \(m\) and
\(n\), conditioning on the number \(k\) chosen from the first
population.

%%%%%%%%%%%%%%%%%%%%%%%%%%%%%%%%%%%%%%%%%%%%%%%%%%%%%%%%%%%%
\subsection{Binomial Inversion}
\label{subsec:binomial-inversion}
%%%%%%%%%%%%%%%%%%%%%%%%%%%%%%%%%%%%%%%%%%%%%%%%%%%%%%%%%%%%

Suppose two sequences satisfy

\[
b_n
=
\sum_{k=0}^{n}
\binom{n}{k}a_k.
\]

Then they can be inverted:

\[
\boxed{
a_n
=
\sum_{k=0}^{n}
(-1)^{n-k}
\binom{n}{k}
b_k.
}
\]

This is called binomial inversion.

It is closely related to inclusion--exclusion and later generalizes to
Möbius inversion on partially ordered sets.

%%%%%%%%%%%%%%%%%%%%%%%%%%%%%%%%%%%%%%%%%%%%%%%%%%%%%%%%%%%%
\subsection{Worked Example: Derangements by Inversion}
\label{subsec:worked-derangement-inversion}
%%%%%%%%%%%%%%%%%%%%%%%%%%%%%%%%%%%%%%%%%%%%%%%%%%%%%%%%%%%%

\begin{workedexamplebox}{Recovering Derangements}

Every permutation may be classified by its set of fixed points.

We have

\[
n!
=
\sum_{k=0}^{n}
\binom{n}{k}D_{n-k}.
\]

Equivalently,

\[
n!
=
\sum_{j=0}^{n}
\binom{n}{j}D_j.
\]

Binomial inversion gives

\[
D_n
=
\sum_{k=0}^{n}
(-1)^{n-k}
\binom{n}{k}k!.
\]

This simplifies to the usual inclusion--exclusion form.

\begin{answerbox}
\[
\boxed{
D_n
=
n!
\sum_{j=0}^{n}
\frac{(-1)^j}{j!}.
}
\]
\end{answerbox}

\end{workedexamplebox}

%%%%%%%%%%%%%%%%%%%%%%%%%%%%%%%%%%%%%%%%%%%%%%%%%%%%%%%%%%%%
\subsection{Young Diagrams}
\label{subsec:young-diagrams}
%%%%%%%%%%%%%%%%%%%%%%%%%%%%%%%%%%%%%%%%%%%%%%%%%%%%%%%%%%%%

A Young diagram is a graphical representation of an integer partition.

For a partition

\[
\lambda
=
(\lambda_1,\ldots,\lambda_k),
\]

draw left-aligned rows of boxes with row lengths

\[
\lambda_1,\ldots,\lambda_k.
\]

Young diagrams become fundamental in algebraic combinatorics,
representation theory, and symmetric functions.

%%%%%%%%%%%%%%%%%%%%%%%%%%%%%%%%%%%%%%%%%%%%%%%%%%%%%%%%%%%%
\subsection{Young Tableaux}
\label{subsec:young-tableaux-preview}
%%%%%%%%%%%%%%%%%%%%%%%%%%%%%%%%%%%%%%%%%%%%%%%%%%%%%%%%%%%%

A Young tableau fills the boxes of a Young diagram with entries.

A standard Young tableau of size \(n\) uses each of

\[
1,2,\ldots,n
\]

exactly once and increases along rows and columns.

The number of standard Young tableaux of a fixed shape is determined by
the hook-length formula.

This topic will be developed more fully in algebraic combinatorics.

%%%%%%%%%%%%%%%%%%%%%%%%%%%%%%%%%%%%%%%%%%%%%%%%%%%%%%%%%%%%
\subsection{Hook-Length Formula Preview}
\label{subsec:hook-length-preview}
%%%%%%%%%%%%%%%%%%%%%%%%%%%%%%%%%%%%%%%%%%%%%%%%%%%%%%%%%%%%

For a partition \(\lambda\vdash n\), the number \(f^\lambda\) of
standard Young tableaux of shape \(\lambda\) is

\[
\boxed{
f^\lambda
=
\frac{n!}{
\prod_{u\in\lambda}h(u)
}.
}
\]

Here,

\[
h(u)
\]

is the hook length of cell \(u\).

The notation

\[
\lambda\vdash n
\]

is read ``lambda partitions \(n\).''

%%%%%%%%%%%%%%%%%%%%%%%%%%%%%%%%%%%%%%%%%%%%%%%%%%%%%%%%%%%%
\subsection{Enumeration by Statistics}
\label{subsec:enumeration-by-statistics}
%%%%%%%%%%%%%%%%%%%%%%%%%%%%%%%%%%%%%%%%%%%%%%%%%%%%%%%%%%%%

Instead of counting all objects equally, one may refine the count by a
statistic.

If

\[
s:\mathcal{C}\to\mathbb{Z}_{\geq0},
\]

define the polynomial

\[
\boxed{
F_n(q)
=
\sum_{c\in\mathcal{C}_n}
q^{s(c)}.
}
\]

The coefficient

\[
[q^k]F_n(q)
\]

counts objects whose statistic equals \(k\).

This framework includes:

\begin{itemize}
    \item inversion-generating polynomials;
    \item descent-generating polynomials;
    \item rank-generating functions;
    \item weight enumerators.
\end{itemize}

%%%%%%%%%%%%%%%%%%%%%%%%%%%%%%%%%%%%%%%%%%%%%%%%%%%%%%%%%%%%
\subsection{\(q\)-Analogues}
\label{subsec:q-analogues}
%%%%%%%%%%%%%%%%%%%%%%%%%%%%%%%%%%%%%%%%%%%%%%%%%%%%%%%%%%%%

A \(q\)-analogue is a polynomial or rational function depending on a
parameter \(q\) that reduces to a familiar counting quantity when

\[
q\to1.
\]

For example, define the \(q\)-integer

\[
\boxed{
[n]_q
=
1+q+q^2+\cdots+q^{n-1}.
}
\]

As \(q\to1\),

\[
[n]_q\to n.
\]

The \(q\)-factorial is

\[
\boxed{
[n]_q!
=
[1]_q[2]_q\cdots[n]_q.
}
\]

These objects become important in algebraic combinatorics.

%%%%%%%%%%%%%%%%%%%%%%%%%%%%%%%%%%%%%%%%%%%%%%%%%%%%%%%%%%%%
\subsection{Gaussian Binomial Coefficients}
\label{subsec:gaussian-binomial-preview}
%%%%%%%%%%%%%%%%%%%%%%%%%%%%%%%%%%%%%%%%%%%%%%%%%%%%%%%%%%%%

The \(q\)-binomial or Gaussian binomial coefficient is

\[
\boxed{
\begin{bmatrix}
n\\
k
\end{bmatrix}_q
=
\frac{
[n]_q!
}{
[k]_q![n-k]_q!
}.
}
\]

When \(q\) is a prime power, this counts \(k\)-dimensional subspaces of

\[
\mathbb{F}_q^n.
\]

Thus a deformation of ordinary binomial coefficients naturally counts
finite vector-space structures.

%%%%%%%%%%%%%%%%%%%%%%%%%%%%%%%%%%%%%%%%%%%%%%%%%%%%%%%%%%%%
\subsection{Asymptotic Enumeration}
\label{subsec:asymptotic-enumeration}
%%%%%%%%%%%%%%%%%%%%%%%%%%%%%%%%%%%%%%%%%%%%%%%%%%%%%%%%%%%%

Exact formulas may become difficult or impossible to obtain.

In such cases, one seeks asymptotic estimates.

If

\[
a_n\sim b_n,
\]

then

\[
\boxed{
\lim_{n\to\infty}
\frac{a_n}{b_n}
=
1.
}
\]

Asymptotic enumeration describes the dominant growth of a combinatorial
class.

%%%%%%%%%%%%%%%%%%%%%%%%%%%%%%%%%%%%%%%%%%%%%%%%%%%%%%%%%%%%
\subsection{Why Asymptotics Matter}
\label{subsec:why-asymptotics-matter}
%%%%%%%%%%%%%%%%%%%%%%%%%%%%%%%%%%%%%%%%%%%%%%%%%%%%%%%%%%%%

Suppose exact formulas are complicated but one wants to compare the
sizes of two families.

Asymptotic formulas can reveal:

\begin{itemize}
    \item exponential growth rates;
    \item polynomial corrections;
    \item typical structure;
    \item limiting probabilities;
    \item algorithmic feasibility.
\end{itemize}

For instance,

\[
n!
\]

and

\[
2^n
\]

both grow rapidly, but factorial growth eventually dominates fixed-base
exponential growth.

%%%%%%%%%%%%%%%%%%%%%%%%%%%%%%%%%%%%%%%%%%%%%%%%%%%%%%%%%%%%
\subsection{Analytic Combinatorics Preview}
\label{subsec:analytic-combinatorics-preview}
%%%%%%%%%%%%%%%%%%%%%%%%%%%%%%%%%%%%%%%%%%%%%%%%%%%%%%%%%%%%

Analytic combinatorics uses complex analysis to extract asymptotic
information from generating functions.

The general philosophy is

\[
\boxed{
\text{combinatorial structure}
\longrightarrow
\text{generating function}
\longrightarrow
\text{singularity behavior}
\longrightarrow
\text{coefficient asymptotics}.
}
\]

A full treatment requires tools from complex analysis and asymptotic
analysis.

%%%%%%%%%%%%%%%%%%%%%%%%%%%%%%%%%%%%%%%%%%%%%%%%%%%%%%%%%%%%
\subsection{Applications}
\label{subsec:enumerative-combinatorics-applications}
%%%%%%%%%%%%%%%%%%%%%%%%%%%%%%%%%%%%%%%%%%%%%%%%%%%%%%%%%%%%

\begin{applicationbox}

\textbf{Algorithms.}
Enumeration measures the size of search spaces and possible data
structures.

\medskip

\textbf{Probability.}
Exact probability calculations often reduce to counting admissible
configurations.

\medskip

\textbf{Algebra.}
Permutations, tableaux, partitions, and \(q\)-analogues appear
throughout representation theory and algebraic combinatorics.

\medskip

\textbf{Number Theory.}
Integer partitions and partition-generating functions play major roles
in additive and modular number theory.

\medskip

\textbf{Geometry.}
Triangulations, lattice paths, tilings, and polyhedral structures create
important counting sequences.

\medskip

\textbf{Graph Theory.}
Enumeration studies labeled and unlabeled graphs, trees, walks, cycles,
and graph colorings.

\medskip

\textbf{Computer Science.}
Recursive structures such as trees, strings, and parse structures are
naturally counted by generating functions.

\medskip

\textbf{Statistical Mechanics.}
Enumeration of configurations contributes to partition functions and
state-counting models.

\end{applicationbox}

%%%%%%%%%%%%%%%%%%%%%%%%%%%%%%%%%%%%%%%%%%%%%%%%%%%%%%%%%%%%
\subsection{Common Errors}
\label{subsec:enumerative-combinatorics-common-errors}
%%%%%%%%%%%%%%%%%%%%%%%%%%%%%%%%%%%%%%%%%%%%%%%%%%%%%%%%%%%%

\subsubsection{Confusing Set and Integer Partitions}

Set partitions divide a set into blocks.

Integer partitions decompose an integer into positive summands.

They are counted by different sequences.

%%%%%%%%%%%%%%%%%%%%%%%%%%%%%%%%%%%%%%%%%%%%%%%%%%%%%%%%%%%%
\subsubsection{Confusing Compositions and Partitions}

Compositions are ordered.

Integer partitions are unordered.

Thus,

\[
3+2
\]

and

\[
2+3
\]

are distinct compositions but the same partition.

%%%%%%%%%%%%%%%%%%%%%%%%%%%%%%%%%%%%%%%%%%%%%%%%%%%%%%%%%%%%
\subsubsection{Confusing Stirling Numbers of the Two Kinds}

The first kind counts permutations by number of cycles.

The second kind counts set partitions by number of blocks.

%%%%%%%%%%%%%%%%%%%%%%%%%%%%%%%%%%%%%%%%%%%%%%%%%%%%%%%%%%%%
\subsubsection{Ignoring Labels}

A formula for labeled objects generally cannot be used unchanged for
unlabeled objects.

Symmetry must be considered.

%%%%%%%%%%%%%%%%%%%%%%%%%%%%%%%%%%%%%%%%%%%%%%%%%%%%%%%%%%%%
\subsubsection{Assuming Equal Counts Imply the Same Objects}

Two combinatorial classes may have the same counting sequence while
consisting of entirely different objects.

A bijection explains equal cardinality without asserting that the
objects themselves are identical.

%%%%%%%%%%%%%%%%%%%%%%%%%%%%%%%%%%%%%%%%%%%%%%%%%%%%%%%%%%%%
\subsubsection{Using a Recurrence Without Correct Initial Values}

The recurrence and initial conditions together define the counting
sequence.

Incorrect base cases produce incorrect enumeration.

%%%%%%%%%%%%%%%%%%%%%%%%%%%%%%%%%%%%%%%%%%%%%%%%%%%%%%%%%%%%
\subsubsection{Ignoring the Statistic Being Counted}

A generating polynomial such as

\[
\sum_{\pi\in S_n}
q^{\operatorname{inv}(\pi)}
\]

does not merely count permutations.

It refines them according to inversion number.

%%%%%%%%%%%%%%%%%%%%%%%%%%%%%%%%%%%%%%%%%%%%%%%%%%%%%%%%%%%%
\subsubsection{Confusing Exact and Asymptotic Equality}

The notation

\[
a_n\sim b_n
\]

does not mean

\[
a_n=b_n.
\]

It means their ratio approaches \(1\).

%%%%%%%%%%%%%%%%%%%%%%%%%%%%%%%%%%%%%%%%%%%%%%%%%%%%%%%%%%%%
\subsection{Exercises}
\label{subsec:enumerative-combinatorics-exercises}
%%%%%%%%%%%%%%%%%%%%%%%%%%%%%%%%%%%%%%%%%%%%%%%%%%%%%%%%%%%%

\begin{exercise}[1]
Compute

\[
\left[
\begin{matrix}
3\\
1
\end{matrix}
\right],
\qquad
\left[
\begin{matrix}
3\\
2
\end{matrix}
\right],
\qquad
\left[
\begin{matrix}
3\\
3
\end{matrix}
\right].
\]
\end{exercise}

\begin{exercise}[1]
Compute

\[
\left\{
\begin{matrix}
4\\
2
\end{matrix}
\right\}.
\]
\end{exercise}

\begin{exercise}[1]
List all set partitions of

\[
\{1,2,3\}.
\]
\end{exercise}

\begin{exercise}[1]
Determine

\[
B_3.
\]
\end{exercise}

\begin{exercise}[1]
List all integer partitions of \(6\).
\end{exercise}

\begin{exercise}[1]
Determine

\[
p(6).
\]
\end{exercise}

\begin{exercise}[1]
List all compositions of \(4\).
\end{exercise}

\begin{exercise}[1]
Verify that there are

\[
2^3=8
\]

compositions of \(4\).
\end{exercise}

\begin{exercise}[1]
Find

\[
C_4.
\]
\end{exercise}

\begin{exercise}[1]
How many triangulations does a convex hexagon have?
\end{exercise}

\begin{exercise}[1]
How many Dyck paths of semilength \(4\) exist?
\end{exercise}

\begin{exercise}[1]
Determine the inversion number of

\[
\pi=4132.
\]
\end{exercise}

\begin{exercise}[1]
Determine the number of descents of

\[
\pi=42513.
\]
\end{exercise}

\begin{exercise}[2]
Use the Stirling recurrence to compute

\[
\left[
\begin{matrix}
5\\
2
\end{matrix}
\right].
\]
\end{exercise}

\begin{exercise}[2]
Use the second-kind Stirling recurrence to compute

\[
\left\{
\begin{matrix}
6\\
3
\end{matrix}
\right\}.
\]
\end{exercise}

\begin{exercise}[2]
Verify

\[
B_4=15.
\]
\end{exercise}

\begin{exercise}[2]
Prove that

\[
B_n
=
\sum_{k=0}^{n}
\left\{
\begin{matrix}
n\\
k
\end{matrix}
\right\}.
\]
\end{exercise}

\begin{exercise}[2]
Show that the number of compositions of \(n\) into exactly \(k\)
positive parts is

\[
\binom{n-1}{k-1}.
\]
\end{exercise}

\begin{exercise}[2]
Use the previous result to prove

\[
\sum_{k=1}^{n}
\binom{n-1}{k-1}
=
2^{n-1}.
\]
\end{exercise}

\begin{exercise}[2]
Draw the Ferrers diagrams of all partitions of \(5\).
\end{exercise}

\begin{exercise}[2]
Find the conjugate of the partition

\[
(5,3,3,1).
\]
\end{exercise}

\begin{exercise}[2]
Explain why conjugation gives a bijection between partitions with at
most \(k\) parts and partitions whose largest part is at most \(k\).
\end{exercise}

\begin{exercise}[2]
Count monotone lattice paths from

\[
(0,0)
\]

to

\[
(5,4).
\]
\end{exercise}

\begin{exercise}[2]
Count permutations of \(6\) elements having exactly \(2\) fixed points.
\end{exercise}

\begin{exercise}[2]
Verify that

\[
n!
=
\sum_{k=0}^{n}
\binom{n}{k}D_{n-k}.
\]
\end{exercise}

\begin{exercise}[2]
Find the Eulerian numbers for \(S_3\) by listing all permutations and
their descents.
\end{exercise}

\begin{exercise}[3]
Prove the recurrence

\[
\left[
\begin{matrix}
n\\
k
\end{matrix}
\right]
=
\left[
\begin{matrix}
n-1\\
k-1
\end{matrix}
\right]
+
(n-1)
\left[
\begin{matrix}
n-1\\
k
\end{matrix}
\right].
\]
\end{exercise}

\begin{exercise}[3]
Prove the recurrence

\[
\left\{
\begin{matrix}
n\\
k
\end{matrix}
\right\}
=
\left\{
\begin{matrix}
n-1\\
k-1
\end{matrix}
\right\}
+
k
\left\{
\begin{matrix}
n-1\\
k
\end{matrix}
\right\}.
\]
\end{exercise}

\begin{exercise}[3]
Derive the inclusion--exclusion formula

\[
\left\{
\begin{matrix}
n\\
k
\end{matrix}
\right\}
=
\frac1{k!}
\sum_{j=0}^{k}
(-1)^{k-j}
\binom{k}{j}j^n.
\]
\end{exercise}

\begin{exercise}[3]
Prove the Bell recurrence

\[
B_{n+1}
=
\sum_{k=0}^{n}
\binom{n}{k}B_k.
\]
\end{exercise}

\begin{exercise}[3]
Derive the Catalan recurrence using binary trees.
\end{exercise}

\begin{exercise}[3]
Give a bijection between balanced parenthesis strings and Dyck paths.
\end{exercise}

\begin{exercise}[3]
Use the reflection principle to derive

\[
C_n
=
\frac1{n+1}
\binom{2n}{n}.
\]
\end{exercise}

\begin{exercise}[3]
Use Stirling's approximation to show

\[
\binom{2n}{n}
\sim
\frac{4^n}{\sqrt{\pi n}}.
\]
\end{exercise}

\begin{exercise}[3]
Use the previous result to derive

\[
C_n
\sim
\frac{4^n}{\sqrt{\pi}n^{3/2}}.
\]
\end{exercise}

\begin{exercise}[3]
Prove by generating functions that partitions into distinct parts and
partitions into odd parts have the same counting function.
\end{exercise}

\begin{exercise}[3]
Derive the inversion-generating polynomial

\[
\prod_{j=1}^{n}
(1+q+\cdots+q^{j-1}).
\]
\end{exercise}

\begin{exercise}[3]
Prove the Eulerian recurrence.
\end{exercise}

\begin{exercise}[3]
Prove binomial inversion.
\end{exercise}

\begin{exercise}[4]
Study the exponential generating function

\[
e^{e^x-1}
\]

for Bell numbers and derive it using labeled combinatorial classes.
\end{exercise}

\begin{exercise}[4]
Investigate the relation between Stirling numbers of the first kind and
falling factorials.
\end{exercise}

\begin{exercise}[4]
Investigate the relation between Stirling numbers of the second kind
and ordinary powers.
\end{exercise}

\begin{exercise}[4]
Study Euler's pentagonal number theorem and its consequences for the
partition function \(p(n)\).
\end{exercise}

\begin{exercise}[4]
Investigate the Hardy--Ramanujan asymptotic formula for \(p(n)\).
\end{exercise}

\begin{exercise}[4]
Study the hook-length formula for standard Young tableaux.
\end{exercise}

\begin{exercise}[4]
Investigate Gaussian binomial coefficients and their interpretation as
counts of subspaces over finite fields.
\end{exercise}

\begin{exercise}[4]
Study \(q\)-analogues of factorials and binomial coefficients.
\end{exercise}

\begin{exercise}[4]
Investigate the Robinson--Schensted correspondence between permutations
and pairs of standard Young tableaux.
\end{exercise}

\begin{exercise}[4]
Study analytic combinatorics and explain how singularities of
generating functions determine asymptotic coefficient growth.
\end{exercise}

%%%%%%%%%%%%%%%%%%%%%%%%%%%%%%%%%%%%%%%%%%%%%%%%%%%%%%%%%%%%
\subsection{Conceptual Summary}
\label{subsec:enumerative-combinatorics-summary}
%%%%%%%%%%%%%%%%%%%%%%%%%%%%%%%%%%%%%%%%%%%%%%%%%%%%%%%%%%%%

Enumerative combinatorics studies counting sequences associated with
structured finite objects.

A combinatorial class produces

\[
\boxed{
c_n=|\mathcal{C}_n|.
}
\]

Permutations are refined by cycles using Stirling numbers of the first
kind:

\[
\boxed{
\left[
\begin{matrix}
n\\
k
\end{matrix}
\right].
}
\]

Set partitions are refined by blocks using Stirling numbers of the
second kind:

\[
\boxed{
\left\{
\begin{matrix}
n\\
k
\end{matrix}
\right\}.
}
\]

Bell numbers count all set partitions:

\[
\boxed{
B_n
=
\sum_{k=0}^{n}
\left\{
\begin{matrix}
n\\
k
\end{matrix}
\right\}.
}
\]

Integer partitions are counted by

\[
\boxed{
p(n),
}
\]

with generating function

\[
\boxed{
\sum_{n\geq0}p(n)x^n
=
\prod_{m\geq1}\frac{1}{1-x^m}.
}
\]

Catalan numbers satisfy

\[
\boxed{
C_n
=
\frac{1}{n+1}\binom{2n}{n}
}
\]

and count many apparently unrelated structures.

Compositions of \(n\) are ordered decompositions and are counted by

\[
\boxed{
2^{n-1}.
}
\]

Permutation statistics lead to refined enumerations such as:

\[
\boxed{
\text{Stirling numbers},
\quad
\text{Eulerian numbers},
\quad
\text{Mahonian numbers}.
}
\]

Asymptotic enumeration studies the growth of counting sequences.

The broader principle is

\[
\boxed{
\text{structure}
\longrightarrow
\text{decomposition}
\longrightarrow
\text{recurrence or generating function}
\longrightarrow
\text{enumeration}.
}
\]

%%%%%%%%%%%%%%%%%%%%%%%%%%%%%%%%%%%%%%%%%%%%%%%%%%%%%%%%%%%%
\subsection{Section Connections}
\label{subsec:enumerative-combinatorics-connections}
%%%%%%%%%%%%%%%%%%%%%%%%%%%%%%%%%%%%%%%%%%%%%%%%%%%%%%%%%%%%

\chapterconnection{
Enumerative Combinatorics combines the counting principles of
Section~\ref{sec:basic-counting}, the recursive methods of
Section~\ref{sec:recurrence-relations}, and the generating-function
methods of Section~\ref{sec:generating-functions}. The next section,
Extremal Combinatorics, changes the emphasis from counting all
structures to determining the largest or smallest structures satisfying
restrictions. Algebraic Combinatorics develops permutation statistics,
Young tableaux, symmetric functions, \(q\)-analogues, and other
algebraic refinements in greater depth. Probabilistic Combinatorics uses
enumeration to analyze random structures. Graph Theory provides major
enumerative families such as trees, graphs, paths, and colorings.
Number Theory connects through integer partitions and arithmetic
generating functions, while Analysis provides asymptotic tools for
large-scale enumeration.
}

%%%%%%%%%%%%%%%%%%%%%%%%%%%%%%%%%%%%%%%%%%%%%%%%%%%%%%%%%%%%
% SECTION SOURCE TRACKING
%%%%%%%%%%%%%%%%%%%%%%%%%%%%%%%%%%%%%%%%%%%%%%%%%%%%%%%%%%%%

%------------------------------------------------------------
% 6.5 Enumerative Combinatorics
%------------------------------------------------------------
%
% Source basis:
%   - Standard enumerative combinatorics.
%   - Standard permutation enumeration.
%   - Standard partition theory.
%   - Standard generating-function methods.
%
% Used for:
%   - combinatorial classes and counting sequences;
%   - labeled and unlabeled enumeration;
%   - permutation cycle enumeration;
%   - Stirling numbers of the first kind;
%   - signed Stirling numbers;
%   - Stirling numbers of the second kind;
%   - Bell numbers;
%   - Bell-number recurrence and EGF;
%   - integer partitions;
%   - Ferrers diagrams;
%   - conjugate partitions;
%   - partitions into distinct parts;
%   - partitions into odd parts;
%   - Catalan numbers;
%   - Dyck paths;
%   - balanced parentheses;
%   - binary trees;
%   - polygon triangulations;
%   - lattice paths;
%   - reflection principle;
%   - central binomial coefficients;
%   - Stirling's approximation;
%   - Catalan asymptotics;
%   - permutation inversions;
%   - Mahonian numbers;
%   - permutation descents;
%   - Eulerian numbers and polynomials;
%   - fixed-point enumeration;
%   - rencontres numbers;
%   - integer compositions;
%   - binomial convolution;
%   - binomial inversion;
%   - Young diagrams and tableaux preview;
%   - hook-length formula preview;
%   - q-analogues;
%   - Gaussian binomial coefficients;
%   - asymptotic enumeration;
%   - analytic combinatorics preview;
%   - applications and exercises.
%
% Notes:
%   - Group-action-based unlabeled enumeration is treated later in
%     Algebraic Combinatorics.
%   - Young tableaux, symmetric functions, and q-analogues are
%     developed more fully in Algebraic Combinatorics.
%   - Integer partition theory connects strongly with Number Theory.
%   - Asymptotic methods connect with Analysis and analytic
%     combinatorics.
%
%%%%%%%%%%%%%%%%%%%%%%%%%%%%%%%%%%%%%%%%%%%%%%%%%%%%%%%%%%%%